% Principal: GDZ PPN236018647, native 20 lower–21 upper, printed 13–14.
\SourceGroup{III}
\HistoricCitation{[Aus dem Tagebuch.]}
\NumberedItem{1}
\par\indent Factorum Summae in Infinito \(=\frac{\pi\pi}{6}\). Sum[ma] Num[erorum]
\begin{flushright}[1796] 20. Jun. G[ottingae].\end{flushright}
\ClosingRule
\NumberedItem{2}
\par\indent Numerus fractionum inaequalium, quarum denominatores certum limitem non superant, ad numerum fractionum omnium, quarum num[eratores] aut denom[inatores] sint diversi infra eundem limitem, in infinito ut \(6:\pi\pi\)
\begin{flushright}[1796] Sept. 6.\end{flushright}
\ClosingRule
\clearpage
\SourceLeaf{21}\ArithmeticHead{14}{L}{ARITHMETIK.\quad NACHLASS.}
\HistoricCitation{[Aus den »Exercitationes Mathematicae«.]}
\NumberedItem{3}
\[
\text{Lim.}\ \frac{\sum\text{ pr. ad }n\text{ infra }n}{nn}
\]
\noindent (Pr. usque ad \(P+30n\)) sunt
\[
(4-1)(3^2-1)(5^2-1)n=576n
\]
\noindent ideo a
\[
P\ .\,.\,.\ Pn\quad\text{generaliter}\quad
\frac{1.3}{2.2}.\frac{2.4}{3.3}.\frac{4.6}{5.5}.\,.\,.\,nn.
\]
\noindent Si \(P:n=\infty:1\)
\par\noindent adeoque
\[
\text{Limes quaesitus}=\frac12.\frac{1.3}{2.2}.\frac{2.4}{3.3}.\frac{4.6}{5.5}.\,.\,.
=\frac3{\pi\pi}=\frac3{9,8696}=0,3039+
\]
